\documentclass{tietreport}

% Import images
\usepackage{graphicx}
\usepackage{booktabs}

% Bibliography configuration
\usepackage[
  date=year,
  style=alphabetic,
  backend=biber,
  sorting=nty,
  maxnames=5,
  minnames=3,
  maxalphanames=4,
  minalphanames=3,
  backref=true,
  doi=false,
  isbn=false,
  url=false,
  eprint=false
]{biblatex}
\addbibresource{references.bib}

%% ==================================================
%% PROJECT FIELDS
%% ==================================================

% Course/Lab header
\TitlePageHeader{%
  UCS503P - Software Engineering (Project)%
}

% Main project title
\ProjectTitle{%
  Window%
}

% Subtitle or project description
\ProjectSubTitle{%
  A Swipe-Based Local Marketplace for Android and the Web%
}

% Student information
\author{%
  \texttt{1024030167} Piyush Malik \\
  \texttt{1024030160} Paarth Singh \\
  \texttt{1024030550} Bhuvik Garg%
}

% Group number, degree, year, program
\TitlePageSubText{%
  Group: \texttt{XXXX} BE Third Year, CSE%
}

% Faculty advisor/guide name
\AdvisorName{%
  Miss Paramveer Kaur%
}

% Evaluation period and department info
\TitlePageFooterText{{%
    \large Mid-Semester Evaluation \\[0.2cm]
    \bfseries Computer Science and Engineering
    Department%
  } \\[0.15cm]
  Thapar Institute of Engineering and Technology,
  Patiala%
}

%% ==================================================

\begin{document}

% Generate title page
\maketitle

% Table of contents (auto-generated from chapters)
\tableofcontents

% ===== CHAPTER 1: INTRODUCTION =====

\chapter{Introduction}

\section{Background and Context}

Small sellers in a city like Patiala, whether they stitch kurtas, sell
juttis, paint, or furnish hostel rooms, rarely have a storefront that
reaches beyond the street they are on. The marketplaces that do reach
customers are built for warehouses: grids of thousands of listings,
mandatory prices, carts, and a checkout that puts a company between the
buyer and the person who made the thing. For a local seller that
distance is the problem, not the solution. The buyer cannot ask whether
a kurta comes in a larger size, the seller cannot say ``come by the shop
this evening'', and neither side knows whether the item on the screen
still exists.

Window is a mobile-first marketplace that keeps the shop and the buyer
within arm's reach of each other. Products are shown one at a time as
cards, like a dating app: a swipe to the left passes, a swipe to the
right tells the seller you are interested, and the seller accepts the
like to open a chat. Prices are a range or nothing at all; the
conversation closes the sale. Everything is local: the feed only shows
sellers within a radius of the buyer's own area, and the radius widens
only when there is not enough nearby.

The project was built by the three of us over the first half of the
semester under the course's constraints: ship a usable increment within
two weeks, iterate weekly, and let continuous integration and delivery
carry every release. The result is a website at
\texttt{window-beige.vercel.app} and a signed Android application
distributed through GitHub Releases, both built from one codebase.

\subsection{Problem Statement}

Buyers and local sellers lack a low-friction, trustworthy channel to
discover each other and talk. Concretely:

\begin{itemize}
  \item \textbf{Decision fatigue.} Catalogue grids present hundreds of
    choices at once; buyers scroll past most of them and remember none.
  \item \textbf{No direct line to the seller.} Questions about sizes,
    materials, delivery, or bargaining have nowhere to go.
  \item \textbf{Stale availability.} A listing says ``in stock'' long
    after the last piece is gone, and the buyer finds out only after
    ordering.
  \item \textbf{No locality.} Platforms show the same products to
    everyone, so a seller two streets away is buried under national
    inventory.
\end{itemize}

Window's hypothesis is that a fast, playful discovery loop (one card at
a time), a direct chat opened by mutual interest, and live stock updates
pushed to everyone who liked a product together turn browsing into
conversations and conversations into local sales.

\section{Project Objectives}

\begin{enumerate}
  \item \textbf{Objective 1 (discovery loop):} Deliver a working swipe
    feed where a right swipe reaches the seller and a seller's
    acceptance opens a chat, with a median swipe-to-match latency of at
    most 2 seconds measured on the client, within the first two-week
    iteration.
  \item \textbf{Objective 2 (local, real-time marketplace):} Show only
    sellers within an adaptive radius (10, 25, then 50 km) of the
    buyer's chosen area, deliver chat messages in under 1 second, and
    propagate a seller's stock change to every interested buyer's
    device within seconds, by the mid-semester evaluation.
  \item \textbf{Objective 3 (seller onboarding and tooling):} Let a
    seller open a shop from their phone in under five minutes (name,
    tags, map location, storefront photo, verified contact) and list a
    product with the camera, a crop step and a price range, with upload
    quotas and simulated purchases, by the mid-semester evaluation.
  \item \textbf{Objective 4 (engineering process):} Keep every change
    behind an automated pipeline: linting, type checking, unit tests,
    security-rule tests and an end-to-end test on a local Firebase
    emulator on every push; automatic production deployment of the
    website; and a signed Android APK built and published for every
    tagged version.
\end{enumerate}

These objectives are specific, measurable through instrumentation
built into the application, achievable with managed services rather
than research, relevant to the course's emphasis on time-to-value, and
bounded to the first half of the semester.

% ===== CHAPTER 2: METHODOLOGY & DESIGN =====

\chapter{Methodology and Design}

\section{System Architecture}

\subsection{High-Level Design}

Window is a single-page web application that is served in two ways:
Vercel hosts it as the website, and Capacitor wraps the identical build
into an Android application. There is no custom server. The client
talks directly to Firebase Authentication and Cloud Firestore, and
Firestore security rules decide who may read and write what. Images are
uploaded from the device to Cloudinary through an unsigned preset; the
database stores only their URLs. Maps come from OpenStreetMap tiles via
Leaflet, with Nominatim for free geocoding.

This ``no backend'' shape was a deliberate trade-off. It let the three
of us ship the first working version in one night and keeps the running
cost at zero on free tiers. The cost is that fan-out logic that a server
would normally perform (for example notifying every buyer who liked a
product when its stock changes) is performed by the seller's client in a
single batched write, and the rules verify that the client is entitled
to each document it touches. Firestore's \texttt{getAfter()} is used so
that a like and the match it creates can be validated together inside
one atomic batch.

The application has two faces that share one account. The shopper face
keeps the middle of the screen free for cards and puts its controls in
the four corners: swipe and super-swipe balances top-left, a switch to
the seller face top-right, liked products and chats bottom-left, and the
profile bottom-right. The seller face has the same switch top-right, a
plus button bottom-centre that opens the camera, shop management
bottom-left and the seller profile bottom-right.

\subsection{Technical Stack}

\begin{itemize}
  \item \textbf{Framework/Language:} React 19 with TypeScript, built by
    Vite; Tailwind CSS 4 for styling; \texttt{motion} for the card
    gestures and the flip, overlay and prompt animations; React Router
    for navigation.
  \item \textbf{Mobile:} Capacitor 8 wraps the web build into Android,
    with native plugins for geolocation, camera and Google Sign-In.
  \item \textbf{Database and authentication:} Cloud Firestore and
    Firebase Authentication (email/password and Google), with security
    rules kept in the repository and deployed by the pipeline.
  \item \textbf{Geospatial:} \texttt{geofire-common} geohashes on every
    product, Leaflet with OpenStreetMap tiles, and Nominatim for forward
    and reverse geocoding.
  \item \textbf{Images:} Cloudinary unsigned uploads after on-device
    compression and cropping (\texttt{react-easy-crop}).
  \item \textbf{Tools:} pnpm, oxlint, Vitest with Testing Library, the
    Firebase emulator suite with \texttt{@firebase/rules-unit-testing},
    GitHub Actions, Vercel, and MkDocs for the documentation site.
\end{itemize}

\section{Implementation Details}

\subsection{Module 1: Discovery Feed and Swipe Economy}

Products are shown as a stack of three cards. The top card can be
dragged, tapped to flip, or driven with the keyboard. Dragging left
lights the left edge of the screen red and stamps the card
\textsc{nope}; dragging right lights the right edge green and stamps it
\textsc{like}. The back of a card shows the price range, the full
description, tags, the shop, and whether the product costs a swipe or a
super swipe.

Every buyer receives five free swipes a day and two super swipes on the
first login from the Android app. A left swipe is free and merely
remembered so the card does not return. A right swipe costs one swipe
and writes, in one batch, the swipe record, a \emph{like} document in a
pending state, and the buyer's new balance. A super swipe costs one
super swipe and claims the product: the like starts accepted and the
match document is created in the same batch, so the chat opens
immediately. Products marked ``super swipes only'' refuse plain swipes,
and the rules enforce this on the server side. When a buyer's balance
reaches zero the counter grows and an overlay offers super swipes,
bundles (simulated payment), or, on the website, the app with its two
free super swipes.

The feed is local. Each product carries a geohash copied from its shop's
pinned location. The client asks Firestore for the geohash ranges that
cover the buyer's radius, filters by exact distance, availability and
swipe history on the device, and ranks products whose tags match the
buyer's chosen interests first and nearest first. If fewer than about
twenty products are found the radius widens from 10 to 25 to 50 km, and
a chip lets the buyer override it.

\subsection{Module 2: Likes, Matches and Chat}

A seller's home screen lists likes waiting for a decision. Accepting
flips the like to accepted and creates the match document; declining
closes it. A match holds denormalised copies of the product, the shop,
both names and unread counters, and if the shop has configured an
automatic message it is copied onto the match so the chat opens with it.
Shops can also buy a day of ``auto-matching'' (99.99 rupees, simulated)
during which every right swipe matches instantly.

Chat is a messages subcollection under the match, read with real-time
listeners. Sending a message is a batch that appends the message and
updates the match's last-message fields and the other party's unread
counter. Availability changes fan out from the seller's client: the
product is updated, every match on that product gets the new stock
state, and every matched buyer receives a notification document, all in
one batch that the rules verify match by match.

\subsection{Module 3: Seller Onboarding and Product Creation}

A seller sets up a shop in four steps shown on a progress rail: name,
up to three tags and an optional website; location on a map with a GPS
button, a place search and a draggable pin; a storefront photo; and the
owner's name, phone and email, each verified by a one-time code (the
delivery of codes is simulated for the prototype and accepts
\texttt{0000}). Google Sign-In fills the owner's name and verified email
in one tap. A shopper who already has an account becomes a seller
without creating a second account; the two faces share one identity.

The plus button opens a live camera preview (or the gallery). The photo
is compressed on the device, cropped to a 3:4 frame, then described:
tags from a curated vocabulary that deliberately excludes food, medical
items, drugs and electronics; a name, for which a product code such as
\texttt{WN-7F3K2Q} is generated on publish; a description limited to 150
words; a price range entered as two numbers that reveal a dual-thumb
slider for fine-tuning; the shop, when the seller owns several; preset
payment modes; and a ``super swipes only'' switch. Listings are metered:
one upload token allows ten products, new sellers receive one, and
bundles of 5, 10, 15 and 20 tokens are sold with the list price struck
through.

\subsection{Module 4: Presentation and Platform}

The application opens on an illustrated storefront with a signboard and
sliding glass doors. On the web, hovering slides the doors ajar; a click
or tap opens them and the camera moves through the doorway into the
location page. On the phone the scene is framed on the door. The tap
that opens the doors is also when the app asks for location permission,
so the map never has to ask unexpectedly.

The design system is a light, storybook style: Nunito display type,
white canvas, one saturated green for headings and primary actions,
blue for links, 12 px radius, 2 px borders and no shadows. Every icon is
an SVG; no emoji are used anywhere.

The pipeline (\texttt{.github/workflows}) runs lint, type check, unit
tests, security-rule tests and an end-to-end test against the Firebase
emulator on every push; builds and uploads a signed release APK on every
push to \texttt{main} and attaches it to the release on every
\texttt{v*} tag; deploys Firestore rules and indexes; and publishes the
documentation site. Vercel deploys production from \texttt{main} and a
preview for every pull request.

% ===== CHAPTER 3: RESULTS & EVALUATION =====

\chapter{Results and Evaluation}

\section{Testing Methodology}

\begin{itemize}
  \item \textbf{Unit Testing:} Vitest covers pure logic: price and time
    formatting, the feed ranking (radius filtering, interest-first
    ordering, distance), product-code generation, word counting, and the
    bundle price tables. 19 tests.
  \item \textbf{Security-rule testing:} \texttt{@firebase/rules-unit-testing}
    exercises the Firestore rules on the emulator as a buyer, a seller
    and a stranger: profile balances cannot go negative, roles are
    immutable, shops and products are owner-only, a plain swipe creates
    only a pending like, a super swipe or an auto-matching shop may
    create the match atomically, only the seller decides a like, only
    participants read a chat, purchases are self-only, and reviews
    require a match. 17 tests.
  \item \textbf{Integration (end-to-end) Testing:} One Vitest suite drives
    the real data layer through the Auth and Firestore emulators: a
    seller opens a shop and lists two products (spending uploads from
    the free token), a shopper sees them within 10 km but not from
    150 km away, swipes, is refused on a super-only product, the seller
    accepts, the chat opens with the auto message and a notification,
    a super swipe claims the super-only product, the auto-matcher makes
    a plain swipe match instantly, a stock change fans out, a review is
    written and edited, and a stranger is denied. 8 steps. The same
    suite was run once against the hosted project to validate the
    deployed rules and indexes.
  \item \textbf{Performance Testing:} The application records its own
    evaluation metrics on the device (swipe-to-match, message delivery,
    availability sync) and shows medians on a hidden metrics page;
    numbers below come from this instrumentation on a mid-range Android
    phone on campus Wi-Fi.
  \item \textbf{User Testing:} The three of us and a handful of
    classmates used the APK and the website against the seeded catalogue
    of four demo shops and sixty products in Patiala. Their feedback
    drove the redesign described in the next section.
\end{itemize}

\section{Performance Metrics}

\begin{table}[h]
\centering
\begin{tabular}{|l|r|r|}
  \hline
  \textbf{Metric} & \textbf{Target} & \textbf{Achieved (median)} \\
  \hline
  Swipe $\rightarrow$ match latency & $\le 2$ s & $\approx 0.4$ s \\
  \hline
  Message delivery (sender to receiver) & $\le 1$ s & $\approx 0.5$ s \\
  \hline
  Availability sync (seller to buyer alert) & seconds & $\approx 1$ s \\
  \hline
  Seller: open a shop & $< 5$ min & $\approx 3$ min \\
  \hline
  CI run (lint, types, unit, rules, e2e) & -- & $\approx 2.5$ min \\
  \hline
  Signed APK build and publish & -- & $\approx 3$ min \\
  \hline
\end{tabular}
\caption{Performance results summary (mid-semester build, v0.4.6)}
\label{tab:performance}
\end{table}

\begin{table}[h]
\centering
\begin{tabular}{|l|r|}
  \hline
  \textbf{Automated check} & \textbf{Count / status} \\
  \hline
  Unit tests & 19 passing \\
  \hline
  Security-rule tests & 17 passing \\
  \hline
  End-to-end steps on the emulator & 8 passing \\
  \hline
  Tagged releases with a signed APK & v0.1.0 to v0.4.6 \\
  \hline
\end{tabular}
\caption{Verification summary}
\label{tab:tests}
\end{table}

\section{Analysis}

\begin{itemize}
  \item \textbf{Objectives met.} All four objectives are met at the
    prototype stage: the discovery loop, the local real-time
    marketplace, seller onboarding with camera listings, and the
    pipeline. Latency targets are met with a wide margin because every
    write is a single Firestore batch and every read is a live listener.
  \item \textbf{Where we exceeded.} The pipeline turned out to be the
    most valuable part of the project. Every one of the fourteen
    releases so far went from commit to a published website and a
    signed APK without a manual step, which made it cheap to act on
    user feedback the same day.
  \item \textbf{Challenges.} Firestore has no geographic query, so
    locality had to be built from geohash range scans and client-side
    filtering. A server-less design pushed fan-out and quota logic into
    the client, which in turn pushed the correctness burden onto
    security rules; \texttt{getAfter()} made atomic like-plus-match
    writes possible, but every rule needed its own test. Profiles
    created before the swipe economy existed produced a negative
    balance on the first swipe and were rejected by the rules; the fix
    was to write absolute balances and backfill missing fields on
    login. On the phone, GPS and camera were blocked when testing over
    plain HTTP, which was solved with an HTTPS development server and
    by asking for permissions at a predictable moment.
  \item \textbf{What changed from the proposal.} The proposal described
    a Tauri desktop application; the target became a phone app plus a
    website. Supabase was replaced by Firestore, Firebase Storage by
    Cloudinary (new Firebase projects need a paid plan for Storage), and
    a dark ``field terminal'' design was replaced by the light storybook
    design after classmates found the dark version dull. Payments and
    one-time codes are simulated, which is stated in the interface.
\end{itemize}

% ===== CHAPTER 4: CONCLUSION =====

\chapter{Conclusion}

\section{Summary of Work}

We set out to make local selling feel like swiping through a shop
window. At mid-semester Window is a working product: a shopper picks an
area and interests, swipes real products from sellers around Patiala,
likes or claims one, and talks to the seller; a seller opens a shop from
a phone, photographs products, accepts likes, chats, and pushes stock
updates. The same code runs as a website and as a signed Android
application, every change passes an automated suite of 44 checks, and
every tagged version is published automatically. All four objectives
were met.

\section{Key Findings}

\begin{itemize}
  \item \textbf{Finding 1:} A marketplace with no server of its own is
    viable when the database's rules are treated as the application's
    contract and tested as rigorously as code; the rule suite caught
    more real bugs than the unit tests did.
  \item \textbf{Finding 2:} Locality is what makes the swipe feel
    meaningful. With a 10 km radius and interest-first ranking, testers
    recognised shops and areas on the cards, which a global feed never
    produced.
  \item \textbf{Finding 3:} The economics layer (free daily swipes,
    super swipes that skip the seller's approval, upload tokens for
    sellers) changed behaviour even though no money moves: testers
    hoarded super swipes for products they actually wanted, which is
    exactly the signal a seller needs.
\end{itemize}

\section{Future Work and Recommendations}

\begin{enumerate}
  \item Real payments through UPI and a real one-time-code provider,
    replacing the simulated purchase and \texttt{0000} flows, with a
    small server or cloud functions to hold the secrets.
  \item Push notifications when the app is closed, which requires the
    paid Firebase plan or a third-party push service.
  \item Moderation and trust: image screening for the tag policy,
    seller verification beyond the self-declared flag, and reporting.
  \item Recommendation beyond ranking by interest and distance, for
    example learning from passes and likes~\cite{firestore}.
  \item An iOS build from the same Capacitor project, and a Play Store
    listing with a production signing key and privacy policy.
  \item Offline resilience: queueing swipes and messages while the
    phone has no signal, which Firestore's local cache makes feasible.
\end{enumerate}

\section{Final Remarks}

The most important lesson for us was the value of shipping small,
verified increments. Because the pipeline made each release cheap, we
could afford to be wrong about the design twice and fix it the same
day, and the product that exists now is the one our testers asked for
rather than the one we first imagined. Window is a marketplace, but it
is also a demonstration that three students with managed services and a
disciplined pipeline can build and operate a real two-sided application
in half a semester.

% ===== BIBLIOGRAPHY =====

% Print bibliography with heading in TOC
\printbibliography[heading=bibintoc]

\end{document}
